\section{What sharper locations would give}\label{sec:frontier}
For the cubic approximation in \eqref{eq:shared-residual}, the first question is how much of its position error is shared between neighboring roots. The same source mean and odd cumulants enter both predictions. A useful estimate would preserve that dependence when taking the difference, and give a uniform error in units of the local spacing. Its comparison is the best bound obtained by differencing or differentiating the existing inverse remainder with the same source data. A separate way to gain height would be to quantify the onset of the inverse as the window constant grows. The complete source must then keep the centered transform away from zero on the chosen complex disk and preserve the count of its inverse.

At one specified all-real \(T\), the product formula supplies a second route to high labels. The finite-order inequalities \eqref{eq:factor-interval} give starting bounds on an individual zero. The question is how the required moment order and precision grow with \(k\) to put two consecutive zeros in intervals narrower than their local spacing. At \(T=1/2\), Section~\ref{sec:comparison} provides a direct comparison with a published counted approximation. Any information from that count used to initialize a product calculation must enter the comparison on both sides.

For a counted family at \(T\), Section~\ref{sec:exterior} makes the accuracy required of a location estimate visible in \(g_k^2 M_k-2\). The leading lattice terms cancel. The error in the complete expression, including zeros outside the located family, must be smaller than the surviving term to establish its sign. A backward estimate from this single time would then have to control
\[
 D_k(T-s)=D_k(T)-\int_0^s
       \bigl(2-g_k(T-u)^2S_{k,1,1}(T-u)\bigr)\dd u
\]
at every partial time, with a common count isolating the same pair. Repeated time differentiation of the mixed moments raises their orders. A usable time remainder must bound those higher orders and all unlocated zeros as the time interval lengthens. The Gaussian evolution of the same upper-time product in \eqref{eq:whole} supplies a direct comparison. A conclusion about all zeros would also require the central pair, the intervening heights and a count excluding nonreal zeros. For the bound in Section~\ref{sec:comparison}, the question is which zero-free region this information establishes at stronger parameters.
